% Module 1 section file. Included by ../module1_stellarator_equilibrium_geometry.tex.
\section{Verification and validation}
\label{sec:verification}

\subsection{Definitions}
\textbf{Verification} asks whether the equations were solved correctly. \textbf{Validation} asks whether those equations and inputs adequately represent the physical system for the intended use. A solver can be verified but invalid for a given experiment, or physically appropriate but implemented incorrectly.

\subsection{Code verification tests}
A Module 1 implementation should include at least the following tests.

\subsubsection{Dimensional and convention tests}
Every exported quantity must carry units or an explicit normalization. Tests should check angle ranges, coordinate orientation, flux signs, and the relation $q=1/\iota$ under the project convention.

\subsubsection{Divergence-free test}
Evaluate
\begin{equation}
\epsilon_{\nabla\cdot B}
=\frac{L_{\mathrm{ref}}\|\nabla\cdot\mathbf B\|_2}{\|\mathbf B\|_2}.
\end{equation}
The norm should decrease with increased numerical resolution until limited by solver tolerance or interpolation error.

\subsubsection{Surface-tangency test}
Evaluate
\begin{equation}
\epsilon_{B\cdot\nabla s}
=\frac{\|\mathbf B\cdot\nabla s\|_2}
{\|\mathbf B\|_\infty\|\nabla s\|_2+\epsilon_{\mathrm{floor}}}.
\end{equation}
A nested-surface equilibrium should make this small throughout the valid domain.

\subsubsection{Force-balance test}
Evaluate
\begin{equation}
\epsilon_{\mathrm{force}}
=\frac{\|\mathbf J\times\mathbf B-\nabla p\|_2}
{\|\nabla p\|_2+\|\mathbf J\times\mathbf B\|_2+\epsilon_{\mathrm{floor}}}.
\end{equation}
The test should also inspect maximum and surface-resolved residuals.

\subsubsection{Field-line transform test}
Trace field lines independently and estimate
\begin{equation}
\iota_{\mathrm{trace}}(s)
=\frac{\Delta\theta}{\Delta\zeta}
\end{equation}
over many turns. Compare with the exported $\iota(s)$. The discrepancy should decrease with tracing tolerance and interpolation resolution.

\subsubsection{Flux conservation test}
Compute magnetic flux through equivalent cross-sections of the same flux tube. Because $\nabla\cdot\mathbf B=0$, the flux should be independent of cross-section up to numerical error.

\subsubsection{Coordinate round-trip test}
For points inside the valid domain, apply
\begin{equation}
(s,\theta,\zeta)\rightarrow\mathbf x
\rightarrow(\tilde s,\tilde\theta,\tilde\zeta).
\end{equation}
Compare $s$ directly and compare angles modulo $2\pi$. Failures near the axis should be treated with coordinate-aware tolerances.

\subsubsection{Manufactured circular-torus test}
Recover the analytic mappings, volume, $V'$, and field-line solution of Section~\ref{sec:analytic_model}. This separates geometry-interface errors from the complexity of a full equilibrium code.

\subsection{Solution verification}
Even a verified code requires case-specific convergence studies. For each equilibrium, report changes in quantities such as
\begin{itemize}[leftmargin=1.6em]
\item magnetic-axis location;
\item $\iota(s)$;
\item volume-averaged beta;
\item minimum Jacobian;
\item selected $B_{mn}$ coefficients;
\item $B_{\min}$ and $B_{\max}$;
\item force residual;
\item derived Module 4 continuum frequencies.
\end{itemize}
The last item is a cross-module sensitivity test: a geometry may appear visually converged while small transform or metric errors still shift resonant frequencies.

\subsection{Validation sources}
Possible validation sources include:
\begin{itemize}[leftmargin=1.6em]
\item comparison with an independently implemented equilibrium code;
\item comparison with published benchmark equilibria;
\item magnetic-diagnostic reconstruction from an experiment;
\item pressure-profile and plasma-boundary measurements;
\item vacuum-field mapping before plasma operation;
\item cross-checks with coil-current perturbation experiments.
\end{itemize}
Validation must be tied to an intended prediction. Agreement in boundary shape does not automatically validate energetic-particle orbit confinement.

\subsection{Uncertainty sources}
Important uncertainties include:
\begin{itemize}[leftmargin=1.6em]
\item coil-position and coil-current errors;
\item plasma-pressure and current-profile uncertainty;
\item diagnostic reconstruction error;
\item finite spectral and radial resolution;
\item coordinate-transformation error;
\item model-form error from scalar-pressure ideal MHD;
\item omitted islands, stochasticity, flow, or anisotropy.
\end{itemize}
Uncertainty should be propagated to quantities actually used by downstream modules, such as $\iota$, orbit-loss fraction, or eigenfrequency.

\section{Equilibrium surrogates and physics constraints}
\label{sec:surrogates}

\subsection{Why build a surrogate?}
A high-fidelity equilibrium solve inside a large parameter scan, control loop, uncertainty analysis, or digital twin can be expensive. A surrogate seeks a rapid approximation
\begin{equation}
\mathcal G:
\left(\text{boundary, coils, profiles, controls}\right)
\mapsto
\left(\mathbf B,\iota,\text{geometry, diagnostics}\right).
\end{equation}
Possible models include neural operators, reduced-basis methods, Gaussian processes, and physics-informed neural networks.

\subsection{Physical constraints}
A useful equilibrium surrogate should be checked against
\begin{align}
\nabla\cdot\mathbf B&=0,\\
\mathbf B\cdot\nabla s&=0,\\
\mathbf J-\frac{1}{\mu_0}\nabla\times\mathbf B&=0,\\
\mathbf J\times\mathbf B-\nabla p&=0,
\end{align}
in the domain where these equations apply. It should also respect periodicity, stellarator symmetry when imposed, boundary conditions, and positive Jacobian orientation.

\subsection{Representation choices}
Predicting Cartesian components of $\mathbf B$ independently can violate $\nabla\cdot\mathbf B=0$. Better choices may include:
\begin{itemize}[leftmargin=1.6em]
\item predicting a vector potential $\mathbf A$ and setting $\mathbf B=\nabla\times\mathbf A$;
\item predicting flux-surface coordinates and using a divergence-free contravariant representation;
\item predicting spectral coefficients constrained by symmetry;
\item learning corrections to a reduced conventional equilibrium rather than the full field.
\end{itemize}
Each representation moves some physical constraints from the loss function into the architecture.

\subsection{Validation requirements}
Low average data error is insufficient. A surrogate should be tested for:
\begin{enumerate}[leftmargin=1.6em]
\item in-distribution interpolation;
\item out-of-distribution boundary and profile changes;
\item force and divergence residuals;
\item topology preservation and positive Jacobian;
\item accuracy of derived quantities such as $\iota$, $k_\parallel$, curvature, and orbit invariants;
\item uncertainty calibration;
\item graceful failure or rejection outside the training domain.
\end{enumerate}

\subsection{Do not replace the reference solver}
The conventional equilibrium solver remains the reference for generating training data, checking residuals, and identifying surrogate failure. The surrogate accelerates selected workflows; it does not eliminate the need for conventional physics verification.

\section{Limitations of the baseline Module 1 model}
\label{sec:limitations}

The textbook description and planned interface are broader than the current executable repository capability. The baseline does not yet implement a validated three-dimensional equilibrium solver. Furthermore, the primary physical model described here has the following limitations unless explicitly extended:
\begin{enumerate}[leftmargin=1.6em]
\item scalar, isotropic pressure rather than an anisotropic pressure tensor;
\item static equilibrium without strong bulk flow;
\item ideal MHD rather than resistive, two-fluid, or kinetic equilibrium;
\item globally nested flux surfaces;
\item no self-consistent magnetic islands or stochastic regions;
\item no scrape-off-layer or open-field-line equilibrium;
\item no direct engineering stress, thermal, or structural coil calculation;
\item no guarantee of MHD stability or good particle confinement;
\item no automatic electrostatic-potential solution;
\item no uncertainty-calibrated experimental reconstruction.
\end{enumerate}
These boundaries are not defects when stated clearly. They define the domain in which Module 1 outputs may be trusted.

\section{Conceptual map for interview-level explanation}
\label{sec:interview_map}

A concise but technically correct explanation is:

\begin{physicsbox}
A stellarator equilibrium is a three-dimensional arrangement of magnetic field, plasma current, and pressure that satisfies $\nabla p=\mathbf J\times\mathbf B$, $\nabla\times\mathbf B=\mu_0\mathbf J$, and $\nabla\cdot\mathbf B=0$. The field is organized into nested toroidal magnetic surfaces. Field lines wind poloidally and toroidally on each surface, with average pitch measured by the rotational transform $\iota$. Surface shape, field-strength variation, metric coefficients, curvature, and $\iota$ determine particle orbits, Alfven-wave propagation, resonances, and transport. A stellarator obtains this geometry primarily from three-dimensional external coils, so equilibrium computation and coordinate transformation are foundational inputs to every later kinetic-MHD module.
\end{physicsbox}

The following distinctions should remain explicit:
\begin{itemize}[leftmargin=1.6em]
\item equilibrium versus stability;
\item vacuum field versus finite-pressure plasma equilibrium;
\item geometric angles versus straight-field-line angles;
\item magnetic surface label versus physical radius;
\item rotational transform $\iota$ versus safety factor $q$;
\item accurate field values versus accurate derivatives;
\item a nested-surface representation versus a field containing islands or stochasticity;
\item a surrogate prediction versus a verified conventional solution.
\end{itemize}

